\documentclass[aip,jcp,reprint,noshowkeys,superscriptaddress]{revtex4-1}
\usepackage{graphicx,dcolumn,bm,xcolor,microtype,multirow,amscd,amsmath,amssymb,amsfonts,physics,longtable,wrapfig,txfonts}
\usepackage[version=4]{mhchem}

\usepackage[utf8]{inputenc}
\usepackage[T1]{fontenc}
\usepackage{txfonts}

\usepackage[
	colorlinks=true,
    citecolor=blue,
    breaklinks=true
	]{hyperref}
\urlstyle{same}

\newcommand{\ie}{\textit{i.e.}}
\newcommand{\eg}{\textit{e.g.}}
\newcommand{\alert}[1]{\textcolor{red}{#1}}
\usepackage[normalem]{ulem}
\newcommand{\titou}[1]{\textcolor{red}{#1}}
\newcommand{\trashPFL}[1]{\textcolor{r\ed}{\sout{#1}}}
\newcommand{\PFL}[1]{\titou{(\underline{\bf PFL}: #1)}}

\newcommand{\cre}[1]{\Hat{a}_{#1}^{\dag}}
\newcommand{\ani}[1]{\Hat{a}_{#1}}
\newcommand{\ca}[2]{\Hat{a}^{#1}_{#2}}

\newcommand{\mc}{\multicolumn}
\newcommand{\fnm}{\footnotemark}
\newcommand{\fnt}{\footnotetext}
\newcommand{\tabc}[1]{\multicolumn{1}{c}{#1}}
\newcommand{\QP}{\textsc{quantum package}}
\newcommand{\T}[1]{#1^{\intercal}}

% coordinates
\newcommand{\br}{\boldsymbol{r}}
\newcommand{\bx}{\boldsymbol{x}}
\newcommand{\bI}{\boldsymbol{1}}
\newcommand{\cK}{\mathcal{K}}

% orbital energies
\newcommand{\eps}{\epsilon}
\newcommand{\ph}{\text{ph}}
\newcommand{\pp}{\text{pp}}
\newcommand{\irr}{\text{irr}}

% shortcuts for greek letters
\newcommand{\si}{\sigma}
\newcommand{\la}{\lambda}
\newcommand{\ii}{\mathrm{i}}
\newcommand{\up}{\uparrow}
\newcommand{\dw}{\downarrow}
\newcommand{\upup}{\up\up}
\newcommand{\updw}{\up\dw}
\newcommand{\dwup}{\dw\up}
\newcommand{\dwdw}{\dw\downarrow}
% addresses
\newcommand{\LCPQ}{Laboratoire de Chimie et Physique Quantiques (UMR 5626), Universit\'e de Toulouse, CNRS, UPS, France}

\begin{document}	

\title{Notes on the parquet and Baym-Kadanoff approximations}

\author{Antoine \surname{Marie}}
	\email{amarie@irsamc.ups-tlse.fr}
	\affiliation{\LCPQ}

\author{Pierre-Fran\c{c}ois \surname{Loos}}
	\email{loos@irsamc.ups-tlse.fr}
	\affiliation{\LCPQ}

\begin{abstract}
Here comes the abstract.
%\bigskip
%\begin{center}
%	\boxed{\includegraphics[width=0.5\linewidth]{TOC}}
%\end{center}
%\bigskip
\end{abstract}

\maketitle

%%%%%%%%%%%%%%%%%%%%%%%%%
\section{Random thoughts}
%%%%%%%%%%%%%%%%%%%%%%%%%

Methods similar to the parquet equations were first introduced in the context of many-body theory by de Dominicis and Martin. \cite{DeDominicis_1964a,DeDominicis_1964b}
The parquet formalism consists of a set of coupled equations, i.e. the Bethe-Salpeter, the parquet, and the Schwinger-Dyson equations, which include both single-particle Green's function and two-particle vertex functions. 
The parquet equations, in principle, provide a framework for self-consistent determination of the self-energy and the vertex corrections.
The parquet formulation preserves, by construction, crossing symmetry and the Pauli exclusion principle.
Crossing symmetry means that a vertex in one channel can also produce a vertex in all other channels by pulling or crossing the vertex legs and multiplying by appropriate constants.
The parquet equations are essentially a generalization of the Bethe-Salpeter equation.

The Baym-Kadanoff approximation \cite{Baym_1961,Baym_1962} (or $\Phi$-derivable approximations) is similar to the parquet approach and can be used to compute one-body correlation functions but does not include self-consistency at the two-body level.
Baym-Kadanoff is primarily focused on the self-energy and satisfies continuity conditions.
However, it produces two-body correlation functions that violate the Pauli principle.
Parquet does enforce this (crossing) symmetry by construction and attaches primary importance on the two-body scattering vertex $\Gamma$, while $\Sigma$ has a subsidiary role and is obtained via integration of the Dyson-Schwinger equation.
The parquet equations guarantee the self-consistent treatment of both one-body and two-body propagators but break conservation laws.
The infinite-order Baym–Kadanoff approximation is also known as ``fluctuation exchange'' (FLEX) approximation. \cite{Bickers_1989}

%%%%%%%%%%%%%%%%%%%%%%%%%
\section{Crossing symmetry}
%%%%%%%%%%%%%%%%%%%%%%%%%

Let us introduce the composite variable $1 \equiv (\br_1,t_1,\sigma_1)$ so that $\sum_{1} = \sum_{\sigma_1} \sum_{\br_1} \int_0^\beta \frac{dt_1}{\beta}$.
The complete ph vertex is
\begin{equation}
	\Gamma = \frac{1}{2} \sum_i \Gamma(12;34) \cre{1} \ani{2} \cre{4} \ani{3}
\end{equation}
with $\Gamma(12;34) = - \Gamma(42,31) = \Gamma(43;21)$ (crossing relation which is extremely difficult to achieve in practice).

The complete pp vertex is 
\begin{equation}
	\Gamma_P = \frac{1}{2} \sum_i \Gamma_P(12;34) \cre{1} \cre{2} \ani{4} \ani{3}
\end{equation}
with $\Gamma(12;34) = - \Gamma_P(14;32)$ and $\Gamma_P(12;34) = - \Gamma_P(21;34) = - \Gamma_P(43)$ (antisymmetry property).
Both the crossing and antisymmetry properties stem from the Pauli exclusion principle.


The full vertex function $F = \Gamma + \Gamma_P$ contains the fully irreducible terms produced by the input $\Lambda$ (which is fully irreducible) and the reducible part from the repeated scattering terms (iterations). 
At the two-body level, there are three different channels for the reducible vertices: one pp channel $\Gamma^\pp$, and two ph (horizontal and vertical) channels, $\Bar{\Gamma}^\ph$ and $\Gamma^\ph$.
For the pp case, reducible diagrams are defined as the ones that can be separated into two pieces by cutting two fermion lines.
ph irreducibility is tricker as two channels exist.


Diagrams irreducible in the horizontal channel (i.e., those which cannot be split by a vertical cut through two lines) contribute to $\Gamma^\ph$, while diagrams irreducible in the vertical channel contribute to $\Bar{\Gamma}^\ph$. 
The quantities $\Gamma^\ph$ and $\Gamma^\pp$ are analogs for ph and pp states of the irreducible self-energy $\Sigma$ for one-body states.

\begin{equation}
\begin{split}
	\frac{1}{2} \sum_i \Gamma^\ph(12;34) \cre{1} \ani{2} \cre{4} \ani{3} 
	& = - \frac{1}{2} \sum_i \Gamma^\ph(12;34) \cre{4} \ani{2} \cre{1} \ani{3} 
	\\
	& = - \frac{1}{2} \sum_i \Bar{\Gamma}^\ph(42;31) \cre{1} \ani{2} \cre{4} \ani{3} 
\end{split}
\end{equation}
hence $\Gamma^\ph(12;34) = -\Bar{\Gamma}^\ph(42;31)$.
The complete vertices $\Gamma$ and $\Gamma_P$ may be recovered from the irreducible vertices by using the Bethe-Salpeter equations:
\begin{gather}
	\Gamma(12;34) = \Gamma^\ph(12;34) + \Gamma(12;56) G^\ph(56;78) \Gamma^\ph(78;34)
	\\
	\Gamma(12;34) = \Bar{\Gamma}^\ph(12;34) + \Gamma(42;56) G^\ph(56;78) \Bar{\Gamma}^\ph(18;37)
	\\
	\Gamma_P(12;34) = \Gamma^\pp(12;34) + \Gamma_P(12;56) G^\pp(56;78) \Gamma^\pp(78;34)	
\end{gather}
with 
\begin{gather}
	G^\ph(12;34) = G(13) G(42)
	\\
	G^\pp(12;34) = -\frac{1}{2} G(13) G(14)
\end{gather}
One must be very careful not to double-count diagrams in the pp sector.

The complete ($\Gamma$ and $\Gamma_P$) and irreducible vertices ($\Gamma^\pp$, $\Gamma^\ph$, and $\Bar{\Gamma}^\ph$) are usually spin adapted (see below).
For example, $\Gamma^\ph$ is written as a sum of its density ($\Gamma^\ph_\text{d}$) and magnetic ($\Gamma^\ph_\text{m}$) components, while $\Gamma^\pp$ is decomposed in its singlet ($\Gamma^\pp_\text{s}$) and triplet ($\Gamma^\pp_\text{t}$) components.
Similar expressions can be found for $\Bar{\Gamma}^\ph$ (crossing relation), and the approximated forms of $\Gamma$ and $\Gamma_P$ (but usually violate crossing conditions).


%%%%%%%%%%%%%%%%%%%%%%%%%
\section{Parquet equations}
%%%%%%%%%%%%%%%%%%%%%%%%%

Let us first focus on the two-body vertices following the presentation of Bickers. \cite{Bickers_2004}
The so-called parquet equations are
\begin{align}
	\begin{split}
    	\Gamma^\ph(12;34) 
    	= \Lambda^\irr(12;34) 
    	& + \Gamma(42;56) G^\ph(56;78)\Bar{\Gamma}^\ph(18;37)
    	\\
    	& + \Gamma_P(41;56) G^\pp(56;78)\Gamma^\pp(78;32)
	\end{split}
	\\
	\begin{split}
    	\Bar{\Gamma}^\ph(12;34) 
    	= \Lambda^\irr(12;34) 
    	& + \Gamma(12;56) G^\ph(56;78)\Gamma^\ph(78;34)
    	\\
    	& - \Gamma_P(14;56) G^\pp(56;78)\Gamma^\pp(78;34)
	\end{split}
	\\
	\begin{split}
    	\Gamma^\pp(12;34) 
    	= \Lambda_P^\irr(12;34) 
    	& - \Gamma(24;45) G^\ph(56;78)\Bar{\Gamma}^\ph(18;37)
    	\\
    	& + \Gamma(14;56) G^\ph(56;78)\Bar{\Gamma}^\ph(28;37)
	\end{split}
\end{align}
Taking advantage of the relation between $\Gamma^\ph$ and $\Bar{\Gamma}^\ph$, one can reduce the problem to two coupled equations:
\begin{align}
	\begin{split}
    	\Gamma^\ph(12;34) 
    	= \Lambda^\irr(12;34) 
    	& - \Gamma(42;56) G^\ph(56;78)\Gamma^\ph(78;31)
    	\\
    	& + \Gamma_P(41;56) G^\pp(56;78)\Gamma^\pp(78;32)
	\end{split}
	\\
	\begin{split}
    	\Gamma^\pp(12;34) 
    	= - \Lambda_P^\irr(12;34) 
    	& + \Gamma(24;56) G^\ph(56;78)\Gamma^\ph(78;31)
    	\\
    	& - \Gamma(14;56) G^\ph(56;78)\Gamma^\ph(78;32)
	\end{split}
\end{align}
$\Gamma$ and $\Gamma_P$ can be also removed using the Bethe-Salpeter equations written above.
To do so, Let us introduce the shortcut notation:
\begin{equation}
	(\Gamma_A \Gamma_B)(12;34) = \Gamma_A(12;56) \Gamma_B(56;34)
\end{equation}
and $\bI(12;34) = \delta_{13} \delta_{24}$.
With these conventions, we have
\begin{gather}
    \Gamma^\ph(12;34) = + \Lambda^\irr(12;34) - \Phi(42;31) + \Psi(41;32)
    \\
    \Gamma^\pp(12;34) = - \Lambda^\irr(14;32) + \Phi(24;31) - \Phi(14;42)
\end{gather}
with
\begin{gather}
    \Phi(12;34) = \qty[ \Gamma^\ph (\bI - G^\ph \Gamma^\ph)^{-1} G^\ph \Gamma^\ph ](12;34)
    \\
    \Psi(12;34) = \qty[ \Gamma^\pp (\bI - G^\pp \Gamma^\pp)^{-1} G^\pp \Gamma^\pp ](12;34)
\end{gather}
With these notations, the crossing relations are
\begin{gather}
	\qty[ \Gamma^\ph + \Phi ](12;34) = - \qty[ \Gamma^\ph + \Phi ](42;31)
	\\
	\qty[ \Gamma^\ph + \Phi ](12;34) = - \qty[ \Gamma^\pp + \Psi ](14;32)
\end{gather}
To ensure crossing symmetry, one must start the self-consistent process with quantities preserving crossing symmetries.
The easiest way of doing this is to use $\Lambda^\irr = v$  (with direct and exchange parts) for a starting point.

Now, we focus on the self-energy which can be written as a non-scattering and a scattering part
\begin{equation}
	\Sigma(11') = \Sigma_1(11') + \Sigma_2(11')
\end{equation}
where, after various manipulations, one gets
\begin{equation}
\begin{split}
	\Sigma_2(11') 
	= \frac{1}{2} \Bigg\{ 
	& - \frac{1}{2} G(76) \qty[ 
	\Lambda^\irr G^\ph v ](17;1'6) 
	\\
	& + G(67) \qty[ \Lambda^\irr_P G^\pp v_P ](17;1'6)
	\Bigg\}
	\\
	& - G(76) \qty[ \Phi G^\ph v ](17;1'6)
	\\
	& + G(67) \qty[ \Psi G^\pp v_P ](17;1'6)
\end{split}
\end{equation}


%%%%%%%%%%%%%%%%%%%%%%%%%
\section{Spin adaptation}
%%%%%%%%%%%%%%%%%%%%%%%%%
Thanks to the spin adaptation, we have the following picture: the electrons interacting through the exchange of four varieties of composite ``bosons'' made of electrons and holes which themselves strongly interact through the exchange of other bosons: density, magnetic, singlet-pair, and triplet-pair fluctuations. 
Thanks to the crossing symmetry of the complete vertices, the parquet equations automatically build in the nonlinear coupling between dressed electron and boson excitations necessary for full consistency.

We now explicitly treat the spin degrees of freedom starting from the ph vertex functions that we decompose in density and magnetic components.
We hence remove the spin variable from the composite variable $1 \equiv (\br_1,t_1)$.
Defining the density and magnetic operators
\begin{gather}
	d(12) = \frac{1}{\sqrt{2}} \qty[ \cre{1\up}\ani{2\up} + \cre{1\dw}\ani{2\dw} ]
	\\
	m(12) = \frac{1}{\sqrt{2}} \qty[ \cre{1\up}\ani{2\up} - \cre{1\dw}\ani{2\dw} ]
\end{gather}
we have
\begin{gather}
	\Gamma_\text{m}^{+1}(12;34) = \Gamma_{\updw,\updw}(12;34) 
	\\
	\Gamma_\text{m}^{-1}(12;34) = \Gamma_{\dwup,\dwup}(12;34)
\end{gather}
and
\begin{gather}
	\Gamma_\text{d}(12;34) = \Gamma_{\upup,\upup}(12;34) + \Gamma_{\upup,\dwdw}(12;34)
	\\
	\Gamma_\text{m}^{0}(12;34) = \Gamma_{\upup,\upup}(12;34) - \Gamma_{\upup,\dwdw}(12;34)
\end{gather}
Hence, $\Gamma_\text{m} \equiv \Gamma_\text{m}^{0} = \Gamma_\text{m}^{\pm1}$.

Now, let us decompose the pp vertex functions into their singlet and triplet components in very much the same way.
Again, the 3 triplet components
\begin{gather}
	\Gamma_\text{t}^{+1}(12;34) = \Gamma_{\upup,\upup}(12;34)
	\\
	\Gamma_\text{t}^{-1}(12;34) = \Gamma_{\dwdw,\dwdw}(12;34)
	\\
	\Gamma_\text{t}^{0}(12;34) = \Gamma_{\updw,\updw}(12;34) - \Gamma_{\updw,\dwup}(12;34)
\end{gather}
are equal and abbreviated as $\Gamma_\text{t}$ while the singlet component reads
\begin{equation}
	\Gamma_\text{s}^{0}(12;34) = \Gamma_{\updw,\updw}(12;34) + \Gamma_{\updw,\dwup}(12;34)
\end{equation}
We recover the fact that the singlet component is symmetric under the exchange of space-time labels, while the triplets are antisymmetric, i.e.
\begin{gather}
	\Gamma_\text{s}(12;34) = + \Gamma_\text{s}(12;43)
	\\
	\Gamma_\text{t}(12;34) = - \Gamma_\text{t}(12;43)
\end{gather}

We are now ready to write down the parquet equations in a spin-adapted basis.
Defining
\begin{gather}
	\Phi_\text{d/m}(12;34) = \qty[ \Gamma_\text{d/m}^\ph (\bI - G^\ph \Gamma_\text{d/m}^\ph)^{-1} G^\ph \Gamma_\text{d/m}^\ph ](12;34)
	\\
	\Psi_\text{s/t}(12;34) = \qty[ \Gamma_\text{s/t}^\pp (\bI - G^\pp \Gamma_\text{s/t}^\pp)^{-1} G^\pp \Gamma_\text{s/t}^\pp ](12;34)
\end{gather}
we get
\begin{align}
	\begin{split}
		\Gamma_\text{d}^\ph(12;34) 
		= \Lambda_\text{d}^\irr(12;34) 
		& - \frac{1}{2} \Phi_\text{d}(42;31) 
		- \frac{3}{2} \Phi_\text{m}(42;31)
		\\
		& + \frac{1}{2} \Psi_\text{s}(41;32) 
		+ \frac{3}{2} \Psi_\text{t}(41;32)
	\end{split}
	\\
	\begin{split}
		\Gamma_\text{m}^\ph(12;34) 
		= \Lambda_\text{m}^\irr(12;34) 
		& - \frac{1}{2} \Phi_\text{d}(42;31) 
		+ \frac{1}{2} \Phi_\text{m}(42;31)
		\\
		& - \frac{1}{2} \Psi_\text{s}(41;32) 
		+ \frac{1}{2} \Psi_\text{t}(41;32)
	\end{split}
	\\
	\begin{split}
		\Gamma_\text{s}^\pp(12;34) 
		= \Lambda_\text{s}^\irr(12;34) 
		& + \frac{1}{2} \Phi_\text{d}(24;31) 
		- \frac{3}{2} \Phi_\text{m}(24;31)
		\\
		& + \frac{1}{2} \Psi_\text{s}(14;32) 
		- \frac{3}{2} \Psi_\text{t}(14;32)
	\end{split}
	\\
	\begin{split}
		\Gamma_\text{t}^\pp(12;34) 
		= \Lambda_\text{t}^\irr(12;34) 
		& + \frac{1}{2} \Phi_\text{d}(24;31) 
		+ \frac{1}{2} \Phi_\text{m}(24;31)
		\\
		& - \frac{1}{2} \Psi_\text{s}(14;32) 
		- \frac{1}{2} \Psi_\text{t}(14;32)
	\end{split}
\end{align}
Check out this nice symmetry!
Now that we have the parquet equations written in a convenient spin-adapted form, the last step is to write down the self-energy via the Dyson-Schwinger equation.
Doing so displays the four varieties of composite bosons mentioned previously:
\begin{equation}
\label{eq:Sig2_sa}
\begin{split}
	\Sigma_2(11') 
	= \frac{1}{2} \Bigg\{ 
	& - \frac{1}{2} G(76) \qty[ 
	\frac{1}{2} \Lambda_\text{d}^\irr G^\ph v_\text{d} + \frac{3}{2} \Lambda_\text{m}^\irr G^\ph v_\text{m} ](17;1'6) 
	\\
	& + G(67) \qty[ \frac{1}{2} \Lambda_\text{s}^\irr G^\pp v_\text{s} + \frac{3}{2} \Lambda_\text{t}^\irr G^\pp v_\text{t} ](17;1'6)
	\Bigg\}
	\\
	& - G(76) \qty[ \frac{1}{2} \Phi_\text{d} G^\ph v_\text{d} + \frac{3}{2} \Phi_\text{m} G^\ph v_\text{m} ](17;1'6)
	\\
	& + G(67) \qty[ \frac{1}{2} \Psi_\text{s} G^\pp v_\text{s} + \frac{3}{2} \Psi_\text{t} G^\pp v_\text{t} ](17;1'6)
\end{split}
\end{equation}
Finally, the crossing relations become
\begin{gather}
	\Gamma_\text{m}(12;34) = - \frac{1}{2} \Gamma_\text{d}(42;31) + \frac{1}{2} \Gamma_\text{m}(42;31)
	\\
	\Gamma_\text{m}(12;34) = - \frac{1}{2} \Gamma_\text{s}(14;32) - \frac{1}{2} \Gamma_\text{t}(14;32)
\end{gather}

%%%%%%%%%%%%%%%%%%%%%%%%%
\section{Fluctuation-exchange approximation}
%%%%%%%%%%%%%%%%%%%%%%%%%
The FLEX approximation is a variant of the Baym-Kadanoff approximation and then focuses on the self-energy itself.
This lack of two-body self-consistency limits the quantitative accuracy of FLEX compared to parquet.
In FLEX, one sets $\Lambda^\irr = \Gamma^\ph = \Gamma^\pp = v$.
Equation \eqref{eq:Sig2_sa} then becomes
\begin{equation}
\begin{split}
	\Sigma_2(11') 
	= \frac{1}{2} \Bigg\{ 
	& - \frac{1}{2} G(76) \qty[ 
	\frac{1}{2} v_\text{d} G^\ph v_\text{d} + \frac{3}{2} v_\text{m} G^\ph v_\text{m} ](17;1'6) 
	\\
	& + G(67) \qty[ \frac{1}{2} v_\text{s} G^\pp v_\text{s} + \frac{3}{2} v_\text{t} G^\pp v_\text{t} ](17;1'6)
	\Bigg\}
	\\
	& - G(76) \Big[ 
	\frac{1}{2} v_\text{d} (\bI - G^\ph v_\text{d})^{-1} (G^\ph v_\text{d})^2 
	\\
	& + \frac{3}{2} v_\text{m} (\bI - G^\ph v_\text{m})^{-1} (G^\ph v_\text{m})^2 
	\Big](17;1'6)
	\\
	& + G(67) \Big[ 
	\frac{1}{2} v_\text{s} (\bI - G^\pp)^{-1} (G^\pp v_\text{s})^2 
	\\
	& + \frac{3}{2} v_\text{t} (\bI - G^\pp)^{-1} (G^\pp v_\text{t})^2
	\Big](17;1'6)
\end{split}
\end{equation}
Because FLEX is a parquet with non-self-consistent two-body vertices, their expressions are a bit different and two flavours exist.
The first flavor includes the single-fluctuation-exchange diagrams while the second flavor includes also the so-called Aslamazov-Larkin diagrams.
When the vertex functions are spin-diagonalized omitting the Aslamazov-Larkin contributions, we recover the parquet equations where one sets $\Lambda^\irr = \Gamma^\ph = \Gamma^\pp = v$.
It has been shown that the Aslamazov-Larkin diagrams deteriorate the accuracy of the two-body vertices obtained via FLEX.

%%%%%%%%%%%%%%%%%%%%%%%%
\acknowledgements{
This project has received funding from the European Research Council (ERC) under the European Union's Horizon 2020 research and innovation programme (Grant agreement No.~863481).}
%%%%%%%%%%%%%%%%%%%%%%%%

%%%%%%%%%%%%%%%%%%%%%%%%%%%%%%%%
%\section*{Data availability statement}
%%%%%%%%%%%%%%%%%%%%%%%%%%%%%%%%
%The data that supports the findings of this study are available within the article.% and its supplementary material.

%%%%%%%%%%%%%%%%%%%%%%%%
\bibliography{parquet-mr-hedin}
%%%%%%%%%%%%%%%%%%%%%%%%

\end{document}

